\chapter{Results}
\label{ch:results}
\section{What changed during development}
The development results explain two methodological choices that materially affect the final comparison. First, the road representation was changed after testing coordinate diffusion on shorter trajectories. In an earlier matched validation study, a free-coordinate model given a planned route achieved 4.58\% whole-trajectory graph consistency. A progress model with the same route conditioning achieved 100\%. Its speed-distribution discrepancy was 0.98 versus 1.39 metres per second. These figures concern a separate 64-point study on 16,779 training windows. They motivate explicit path decoding. They do not form a numerically comparable baseline row in the 512-point held-out experiment.

Second, candidate paths were diversified before final locking. Eight motion draws on a common route varied principally in their progress along that route. Independently drawing the paths increased the mean number of distinct paths to 6.70 per eight-candidate validation pool. Mean pairwise coordinate separation increased from about 62--67 metres to about 147 metres. The changed opportunity set mattered to interpretation. A classifier restricted to one corridor cannot select an alternative branch, regardless of its recognition accuracy. Yet giving it several branches did not resolve the disagreement between the image and raw-trajectory auditors.

On the 33 validation requests with diverse paths, GAF ranking increased image-auditor macro accuracy from 37.22\% to 63.33\% relative to random selection. Coordinate-auditor accuracy fell from 76.94\% to 72.78\%. The same validation pool had zero planning failures or path-end caps. This motivated testing empirical transfer and downstream usefulness before attributing an advantage to the feedback mechanism.

\section{Held-out coverage and road validity}
All 53 qualified requests completed eight path slots. Each of the three diffusion seeds generated 424 trajectories and each of the three empirical resamplings generated the same number. The diffusion export therefore contains 1,272 candidates comprising 651,264 timed positions. These are repeated conditional draws for 53 held-out participants. Table~\ref{tab:cohorts} gives the population retained at the preceding stages.

All candidates in both families passed the directed graph audit, all could be encoded and none reached the path-end cap. The final masks needed no fallback. The common result has a straightforward interpretation: both motion models operate through a shared permitted-path decoder on the selected population. A geometric advantage of one over the other cannot be identified from two identical pass rates. The coverage result also distinguishes the successful generation of complete pools from the eligibility of only 53 of the original 270 IDs.

\fig{map_examples.pdf}{.94}{First served condition of each supported class at seed 42. Each row compares three ways to select from the same eight blue candidates. Gray lanes are permitted for that class, green marks the given start, and orange marks the selected trace. Each panel uses local SUMO metres.}{map}
Figure~\ref{fig:map} uses a fixed rule established before viewing results. In the bus example random and GAF selection chose the same candidate. For the motorcycle, each selector chose a different candidate. For the private car, random and GAF again coincide. The map shows the spatial consequences of a selection and the routes available to it. It cannot by itself tell whether a bus accelerates realistically. Appendix~\ref{app:cases} plots the timing of these same examples and gives class-level distributions using all eligible participants.

\section{Learning motion beyond empirical transfer}
Table~\ref{tab:w1} compares the complete candidate populations. Every method has access to identical geometry. The learned model improves the point estimate for speed and travel distance, while the empirical transfer baseline has the smaller acceleration discrepancy. Only the travel comparison has a paired interval entirely above zero. It is reasonable to say that diffusion brought the travel-distance distribution closer on the strict test population, conditional on the recorded models and paths.

\tab{Final all-candidate comparisons. Distances are equally weighted across the five supported classes. Positive paired differences favor diffusion.}{w1}{\begin{tabularx}{\linewidth}{Xrrr}\toprule Quantity & Diffusion & Empirical & Difference [95\% interval]\\\midrule
Speed (m/s) &1.244&1.446& $0.203$ [$-0.018$, $0.993$]\\
Acceleration (m/s$^2$) &0.269&0.239& $-0.030$ [$-0.081$, $0.046$]\\
Travel (m) &177.414&201.119& $23.706$ [$7.857$, $132.912$]\\\bottomrule\end{tabularx}}
\fig{paired_w1.pdf}{1}{Empirical-minus-diffusion differences in the final test. The three panels use different physical units. Bars are vehicle-level percentile intervals conditional on fixed models, routes and candidate draws.}{w1}

It is informative to compare directions across development and test, without pooling the populations. On validation the empirical-minus-diffusion acceleration difference was $0.0373$ metres per second squared, interval $[0.0017,0.0799]$. On test the direction reversed and the interval spans zero. The validation travel difference was $6.99$ metres with an interval spanning zero, whereas the test travel interval excludes zero. These changes underscore why selecting the most favorable validation metric and reporting it alone would have been misleading.

The travel result also requires a distinction between absolute and distributional accuracy. It does not imply that the generated car for a particular request traveled 23.7 metres closer to that participant's observed trip. The measured quantity compares distributions of total travel. The wide upper endpoint reflects substantial uncertainty from a small and uneven set of reference vehicles. The class-specific pages in the appendix retain the support behind the macro average.

\subsection{Which classes account for the travel result?}
Averaging the class-specific travel distances in Appendix~\ref{app:cases} recovers the aggregate result in Table~\ref{tab:w1}. Relative to empirical transfer, diffusion reduces travel discrepancy by about 85.5 metres for buses, 20.1 metres for motorcycles, 73.1 metres for private cars and 18.7 metres for bicycles. For taxis it increases discrepancy by 79.0 metres. These are descriptive components of the prespecified macro average, not separate statistically significant class effects.

The finding for cars is supported by 33 reference conditions and that for buses by seven. The unfavorable taxi contribution receives the same weight despite its two references. It accounts for much of the difference between stating that the overall distribution improved and claiming improvement for every class. Giving every candidate equal weight across the entire test set would instead give private cars much greater influence. These would be different estimands. We use the stated macro measure consistently so the class-control question remains visible.

The stop fractions offer a second interpretation of the remaining discrepancy. For motorcycles, the observed set spends 24.6\% of sampled intervals below the stop threshold, while diffusion spends about 7.0\% and transfer about 7.4\%. For taxis the corresponding figures are 19.2\%, 2.2\% and 5.3\%. Both motion mechanisms underrepresent low-speed intervals for these particular references. The current model can generate zero increments, so their mere possibility under the decoder is insufficient to match their observed frequency. Whether this gap persists on a larger independent cohort remains open.


\section{Diversity and the copying control}
Diffusion candidates had mean within-request distances of 137.0, 140.6 and 139.1 metres across seeds. The empirical pools had corresponding distances of 200.6, 206.9 and 203.5 metres. More separation did not imply uniformly closer motion distributions. Indeed, the empirical model had higher travel discrepancy even though its candidate paths spread farther apart at corresponding times. Dispersion and fidelity answer different questions.

The empirical comparator copied a complete training profile by construction. The aligned increment diagnostic identified 100\% of empirical profiles as near copies and 0\% of diffusion profiles. Median nearest-neighbor error was 0.306--0.331 metres per increment for diffusion. The observed held-out reference median was 0.383 metres. The result provides evidence against whole-profile transfer by the diffusion generator at the specified tolerance. A generator could still repeat shorter fragments, time-shifted profiles or common movements with small perturbations. This diagnostic cannot support a privacy claim.

\section{Recognition and the effect of selection}
On all 182 encodable test windows, the self-trained GAF family achieved mean accuracy of 64.29\% and macro-F1 of 53.38\%. The corresponding coordinate TCN achieved 78.57\% and 72.92\%. These are observed-window recognition results, not measures of synthetic trajectory quality. The gap matters when comparing auditors. Each has different errors and access to different features. No auditor is a ground-truth labeler for generated movement.

\tab{Equal-budget selection on the final diffusion pools. Audit columns are class-macro agreement percentages with the requested label. They are distinct from six-class recognition F1.}{select}{\begin{tabular}{lrrrr}\toprule Selector & GAF audit & TCN audit & Speed $\W$ & Travel $\W$\\\midrule
Random &51.85&89.86&1.402&206.00\\
GAF &70.45&78.10&1.410&191.52\\
TCN &47.32&93.18&1.444&209.30\\
Shuffled-label GAF &39.55&76.95&1.501&195.80\\\bottomrule\end{tabular}}
When the selector optimizes image-based class scores, agreement with a separate image auditor rises by 18.60 percentage points over random selection. Agreement with the coordinate auditor drops by 11.76 points. Ranking with a coordinate scorer raises coordinate agreement to 93.18\% but gives a lower image agreement of 47.32\%. These are paired-pool observations, so the reversal cannot be explained by giving one method more samples. It is consistent with selection favoring the features used within each family. The experiments do not identify the individual feature responsible.

There is some metric-specific variation beyond classifier agreement. In the diffusion pools, GAF selection has a smaller travel discrepancy than random selection, 191.52 versus 206.00 metres, with little change in speed discrepancy. Applying the same selector to empirical motions increases travel discrepancy from 233.89 to 252.55 metres. No predeclared paired interval establishes a travel improvement from GAF selection. These comparisons are descriptive and give no basis for transferring the apparent improvement from one generator to another.

\section{Using synthetic samples for training}
\tab{Six-class macro-F1 on held-out observed windows. Each cell averages three final learner seeds. Added examples use training IDs only and each augmented group has the same class inventory.}{utility}{\begin{tabular}{lrr}\toprule Training group & TCN (\%) & GRU (\%)\\\midrule
Observed only &76.92&64.38\\
Repeat observed &72.77&59.70\\
Time warp &73.32&63.14\\
Random synthetic &70.29&60.82\\
GAF-selected synthetic &71.69&62.68\\
TCN-selected synthetic &72.08&61.20\\\bottomrule\end{tabular}}
\fig{utility.pdf}{1}{Change in macro-F1 relative to training on observed data alone. All learners have equal training budgets. The interval accounts for paired resampling of observed test vehicles; the augmentation generator is fixed.}{utility}
The TCN lost 5.23 percentage points when trained with GAF-selected synthetic data, with a paired interval of $[-9.01,-1.65]$. The GRU lost 1.70 points, with interval $[-6.42,3.13]$. Compared with random synthetic selection, GAF selection had positive point differences of 1.40 and 1.86 points, but both intervals include zero. There is insufficient evidence that GAF is a better augmentation rule than random selection for these learners.

The GRU performed better with GAF-selected traces than with repeated source vehicles by 2.98 points, with interval $[0.28,5.95]$. This single contrast is compatible with the observed-only and time-warp groups still doing better in point estimates. Choosing repeated source vehicles as the sole baseline would make the usefulness claim depend on an unusually weak comparator. The six-way table preserves this distinction.

A practical reading is that topology-valid, class-labeled synthetic trajectories can be produced and written back in the input schema. At this augmentation ratio, the selected examples do not repay the opportunity cost of diverting a fixed training budget from observed trajectories. Several simpler augmented groups also perform worse than observed-only training. The data therefore leave open how much of the change is due to motion fidelity, class reweighting, extra correlated source vehicle examples or the reduction in expected presentations per original window.

\subsection{What the class-level errors reveal}
The standalone image classifier's mean observed-test F1 ranges from 85.3\% for pedestrians to 0.0\% for taxis. The coordinate-based family reaches 27.8\% on taxis, still a weak result. Across the three diffusion seeds the mean number of distinct images in an eight-candidate test pool is 8.00, 7.98 and 8.00. Thus the scorer typically sees distinct inputs, even though distinct images need not give useful distinctions among the requested classes.

A second view is to inspect downstream recall. For the TCN, adding GAF-selected examples increases bus recall from 84.6\% to 92.3\%, while motorcycle recall decreases from 51.6\% to 43.0\% and taxi recall from 62.5\% to 45.8\%. For the GRU, bus recall rises from 41.0\% to 59.0\%, while motorcycle recall decreases from 37.6\% to 30.1\%. A change in the frequency with which a learner predicts a class can affect its precision simultaneously. The appendix retains all classes and all groups, so these individual changes are not used to reinterpret the overall macro-F1 effect.

The validation results also explain why the final test was necessary. During development the TCN GAF-versus-observed-only difference was $-1.33$ points and the GRU difference was $+0.90$ points. Both paired intervals spanned zero. On the final cohort the TCN decline was larger and the GRU point estimate changed sign. A claim of successful augmentation based solely on the positive validation GRU estimate would therefore have failed even as a point prediction of the final comparison.

Under uniform sampling, each observed training window has expected exposure $73{,}600/1{,}437$ in the observed-only group and $73{,}600/1{,}862$ in an augmented group. These are roughly 51.2 and 39.5 presentations, respectively. Thus a fixed training budget involves a concrete exchange: additional examples replace a fraction of the presentations of observed windows. The observed-only result should be interpreted at that budget. Holding original-example exposure constant would examine a different perturbation and would require a new experimental protocol.

